\chapter{Modelo}
\label{chap:modelo}

\section{Juego secuencial con $m > 1$ líderes}

El modelo que motiva la estimación se basa en un esquema de competencia secuencial en cantidades (Stackelberg) aplicado al mercado internacional del café, donde los países exportadores interactúan estratégicamente en su toma de decisiones de producción. Esta estructura captura la asimetría entre naciones con diferencias notables de participación de mercado, generando dinámicas de dependencia estratégica en la competencia.

El grano verde de café es considerado un bien homogéneo, a pesar de las diferencias entre especies (principalmente Arábica y Robusta) y categorías comerciales (café de altura o estándar), en base a lo desarrollado en la sección \ref{chap:mercado}.

Ahora bien, un mercado de competencia en cantidades puede tener diferentes aproximaciones. En este caso, un supuesto importante para la estimación es imponer al mercado una estructura de competencia Stackelberg, tanto por la relevancia histórica de algunos competidores en las exportaciones mundiales como por su injerencia pasada en la intensidad competitiva del mercado.

Por otro lado, la elección del modelo de competencia secuencial tiene importantes implicancias metodológicas. Por ejemplo, en el caso de una competencia \textit{à la Cournot}, donde todos los actores deciden sus cantidades de forma simultánea, las funciones de reacción tenderían a homogeneizar artificialmente el comportamiento de los países, independientemente de su tamaño o influencia efectiva.

El siguiente modelo sigue a Huck et al. \citeyearpar{huck2001}. Se define con $n$ a la cantidad de países exportadores en el mercado internacional del café, los que enfrentan una función de demanda lineal $P(Q)=A-B\cdot Q$, donde $Q=\sum_{i=1}^{n}q_{i}$ representa la exportación total y $q_{i}$ es la exportación de un país, con $i \in \{1, ..., n\}$. A su vez, se supone la existencia de $m$ países líderes ($m<n$) y $n-m$ países seguidores. Por último, se consideran costos marginales constantes y únicos para cada país ($c_{i}$), los que dependen de factores de oferta agrícola y condiciones climáticas adversas.

Sean $\bar{Q}^{l}$ la cantidad exportada por los $m$ países líderes y $Q^{s}_{-i}$ la cantidad de $n-m-1$ países seguidores, el seguidor modelo $i$ debe escoger un $q^{s}_{i}$ que maximice:

\begin{equation}
    max_{q_{i}^{s} \ge 0}: \pi^{s}(q_{i}^{s}) = (A-B(\bar{Q}^{l}+Q^{s}_{-i} + q^{s}_{i})- c_{i})\cdot q_{i}^{s} 
\end{equation}

Por otro lado, sea $Q^{l}_{-i}$ la cantidad exportada por los países líderes diferentes al líder modelo $i$, éste debe escoger un $q^{l}_{i}$ que maximice:

$$ max_{q_{i}^{l} \ge 0}: \pi^{l}(q_{i}^{l}) = (A-B(Q^{l}_{-i}+Q^{s}(q^{l}_{i} + Q^{l}_{-i}) + q^{l}_{i})- c_{i})\cdot q_{i}^{l} $$

En equilibrio, cada seguidor resuelve (1) tomando las exportaciones del grupo de líderes como dadas, llevando a su cantidad óptima:

$$ \frac{\partial \pi^{s}}{\partial q_{i}^{s}} = 0 \iff P-c_{i}^{s}-B q_{i}^{s}=0$$
$$ q_{i}^{s}=\frac{P-c_{i}^{s}}{B} = \frac{A-B \cdot \bar{Q}^{l}-B\cdot Q^{s}_{-i}-c_{i}}{2B}$$

Luego, bajo un supuesto de simetría entre los seguidores, es posible agregar el resultado, de manera que $Q^{s}_{-i} = (n-m-1)\cdot q^{s}_{i}$. Así:

\begin{equation}
    q^{s}_{i}(\bar{Q}^{l}) = \frac{A-B\cdot \bar{Q}^{l}-c_{i}}{(n-m+1)\cdot B}
\end{equation}

Por su parte, los líderes resuelven el problema de optimización por inducción hacia atrás, incorporando la mejor respuesta de los seguidores:

$$  max_{q_{i}^{l} \ge 0}: \pi^{l}(q_{i}^{l}) = (A-B(Q^{l}_{-i}+ \textcolor{blue}{\left(\frac{n-m}{(n-m+1)\cdot B}\right)\cdot (A-B(q^{l}_{i}+Q^{l}_{-i})-c_{i})} + q^{l}_{i})- c_{i})\cdot q_{i}^{l} $$

$$ \frac{\partial \pi^{l}}{\partial q^{l}_{i}} = 0 \iff q^{l}_{i} = \frac{A-B\cdot Q^{l}_{-i} - c_{i}}{2B}$$

Nuevamente, bajo un supuesto de simetría, es directo que $Q^{l}_{-i} = (m-1)\cdot q^{l}_{i}$. Por tanto:

$$ q^{l}_{i} = \frac{A-B\cdot (m-1) \cdot q^{l}_{i} - c_{i}}{2B} $$

\begin{equation}
    q^{l}_{i} = \frac{A-c_{i}}{(m+1)\cdot B}    
\end{equation}

Así, de sustituir (3) en (2), se obtiene que:

\begin{equation}
    q^{s}_{i} = \frac{A-c_{i}}{(m+1)(n-m+1)B}
\end{equation}

Es importante notar que las expresiones anteriores se obtienen bajo un supuesto de simetría dentro de cada grupo (líderes y seguidores), lo que permite derivar formas cerradas para ilustrar la intuición económica fundamental del modelo: los seguidores enfrentan una penalización adicional $(n-m+1)$ en sus cantidades óptimas respecto de los líderes.

Sin embargo, en la estimación empírica se relaja este supuesto de simetría. Dado que los costos marginales varían sustancialmente entre países --incluso dentro del mismo grupo-- la implementación empírica resuelve el sistema completo de condiciones de primer orden sin imponer igualdad entre los miembros de cada grupo exportador. Esto permite capturar la heterogeneidad observada en las capacidades productivas de los distintos actores.

Es directo ver cómo las funciones de reacción difieren entre países líderes y seguidores. En el caso de estos últimos, las cantidades óptimas que resuelven la interacción estratégica se ven penalizadas por el factor $(n-m+1)$, es decir, la brecha aumenta cuanto menor sea la cantidad de líderes, o cuanto más se incremente la diferencia $n-m$. 

Reordenando las condiciones de primer orden, es posible obtener expresiones implícitas para los costos marginales, que relacionan precios observados con cantidades individuales. Estas expresiones no requieren asumir simetría y constituyen la base de la estimación estructural.

Para un país seguidor $i$, se tiene que:

\begin{equation}
    \label{eq:c_seguidor}
    c_{i}^{s} = P + \frac{\partial P}{\partial q_{i}^{s}}\cdot q_{i}^{s}
\end{equation}

Por su parte, para un país líder $i$, la condición incorpora el efecto estratégico sobre los seguidores. En el caso de $m$ líderes simétricos, este factor sería igual a $\frac{1}{m+1}$, pero en la estimación con líderes heterogéneos se mantiene la estructura cualitativa:

\begin{equation}
    \label{eq:c_lider}
    c_{i}^{l} = P + \left( \frac{\partial P}{\partial Q^{s}}\right) \cdot \frac{\partial Q^{s}}{\partial q_{i}^{l}}\cdot q_{i}^{l}
\end{equation}

\section{Implementación empírica}

El proceso de estimación procede en tres etapas, relajando el supuesto de simetría del modelo teórico para capturar la heterogeneidad observada.

En la primera etapa, se estima la demanda agregada mediante variables instrumentales, utilizando variables climáticas como instrumentos para la cantidad exportada. Esto permite recuperar la elasticidad-precio de la demanda y construir los costos marginales implícitos país-específicos mediante las ecuaciones \eqref{eq:c_seguidor} y \eqref{eq:c_lider}.

Más adelante, en una segunda etapa, los costos marginales recuperados se modelan como función de variables observables —fertilizantes, tendencias temporales, exposición a estrés térmico—  mediante un modelo de efectos fijos por país. Esta especificación permite cuantificar cómo las variaciones en condiciones climáticas afectan la competitividad relativa de cada productor.

Por último, para evaluar escenarios contrafactuales, se resuelve numéricamente el sistema completo de condiciones de primer orden para todos los países. Formalmente, el equilibrio se caracteriza por el vector de cantidades $(q_1^*, ..., q_n^*)$ que satisface simultáneamente:

\begin{itemize}
    \item Para cada líder $i \in L$: 
    $$P(Q^*) + \frac{\partial P}{\partial q_i}\left(q_i^* + Q_s^*(Q_L^*)\right) 
    \cdot q_i^* \left(1 + \frac{\partial Q_s}{\partial Q_L}\right) = c_i$$
    
    \item Para cada seguidor $j \in S$: 
    $$P(Q^*) + \frac{\partial P}{\partial q_j} \cdot q_j^* = c_j$$
    
    \item Equilibrio de mercado: 
    $$Q^* = \sum_{i \in L} q_i^* + \sum_{j \in S} q_j^*$$
\end{itemize}

donde $Q_s^*(Q_L^*)$ representa la respuesta agregada de los seguidores ante las cantidades de los líderes, y $c_i$, $c_j$ son los costos marginales predichos según las condiciones climáticas del escenario evaluado.

Este sistema se resuelve mediante el algoritmo de Newton-Raphson multivariado, lo que permite obtener el equilibrio sin imponer restricciones de simetría. La ventaja de este enfoque es que preserva la estructura jerárquica del modelo Stackelberg —donde los líderes internalizan la respuesta de los seguidores— mientras captura las diferencias sustanciales en costos y escalas entre países.

Los escenarios contrafactuales se construyen modificando las variables climáticas en las funciones de costo estimadas y resolviendo nuevamente el sistema. Esto permite cuantificar tanto los efectos directos del cambio climático sobre cada productor como los efectos indirectos derivados de la recomposición del equilibrio de mercado.